% ------------------------- Capítulo I -----------------------------
% Dfinições
\newtheorem{def:base-prob}{Definição}[chapter]
\newtheorem{def:pontualmente continua}[def:base-prob]{Definição}
\newtheorem{def:uniformemente continua}[def:base-prob]{Definição}
\newtheorem{def:absolutamente continua}[def:base-prob]{Definição}
\newtheorem{def:medida nula}[def:base-prob]{Definição}
\newtheorem{def:mensuravel}[def:base-prob]{Definição}
\newtheorem{def:medida}[def:base-prob]{Definição}
\newtheorem{def:medida de contagem}[def:base-prob]{Definição}
\newtheorem{def:medida de lebesgue}[def:base-prob]{Definição}
\newtheorem{def:medida completa}[def:base-prob]{Definição}
\newtheorem{def:medida absolutamente continua}[def:base-prob]{Definição}
\newtheorem{def:funcao mensuravel}[def:base-prob]{Definição}
\newtheorem{def:steltjes}[def:base-prob]{Definição}
\newtheorem{def:funcao simples}[def:base-prob]{Definição}
\newtheorem{def:integral de lebesgue}[def:base-prob]{Definição}
\newtheorem{def:lp}[def:base-prob]{Definição}
\newtheorem{def:quase sempre}[def:base-prob]{Definição}
\newtheorem{def:espaco de probabilidade}[def:base-prob]{Definição}
\newtheorem{def:sigma algebra}[def:base-prob]{Definição}
\newtheorem{def:borel}[def:base-prob]{Definição}
\newtheorem{def:medida imagem}[def:base-prob]{Definição}
\newtheorem{def:probabilidade}[def:base-prob]{Definição}
\newtheorem{def:evento equiprovavel}[def:base-prob]{Definição}
\newtheorem{def:eventos encaixados}[def:base-prob]{Definição}
\newtheorem{def:probabilidade condicional}[def:base-prob]{Definição}
\newtheorem{def:independencia}[def:base-prob]{Definição}
% Definições
\newtheorem{exp:base-prob}{Exemplo}[chapter]
\newtheorem{exp:pontualmente contínua}[exp:base-prob]{Exemplo}
\newtheorem{exp:pontualmente contínua neg}[exp:base-prob]{Exemplo}
\newtheorem{exp:uniformemente contínua}[exp:base-prob]{Exemplo}
\newtheorem{exp:uniformemente contínua neg}[exp:base-prob]{Exemplo}
% Proposições
\newtheorem{prop:base-prob}{Proposição}[chapter]
\newtheorem{prop:base calculo probabilidades}[prop:base-prob]{Proposição}
\newtheorem{prop:base calculo probabilidades 2}[prop:base-prob]{Proposição}
\newtheorem{prop:eventos encaixados}[prop:base-prob]{Proposição}
\newtheorem{prop:probabilidade condicional}[prop:base-prob]{Proposição}
\newtheorem{prop:regra da multiplicacao}[prop:base-prob]{Proposição}
\newtheorem{prop:independencia}[prop:base-prob]{Proposição}
\newtheorem{prop:independencia complementar}[prop:base-prob]{Proposição}
% Lema
\newtheorem{lem:base-prob}{Lema}[chapter]
\newtheorem{lem:fatou}[lem:base-prob]{Lema}
% Teorema
\newtheorem{thm:base-prob}{Teorema}[chapter]
\newtheorem{thm:heine}[thm:base-prob]{Teorema}
\newtheorem{thm:mensuravel real}[thm:base-prob]{Teorema}
\newtheorem{thm:calculo medida}[thm:base-prob]{Teorema}
\newtheorem{thm:medida de subconjunto}[thm:base-prob]{Teorema}
\newtheorem{thm:lebesgue sigma finita}[thm:base-prob]{Teorema}
\newtheorem{thm:funcoes mensuraveis}[thm:base-prob]{Teorema}
\newtheorem{thm:medida steltjes}[thm:base-prob]{Teorema}
\newtheorem{thm:aproximacao monotona}[thm:base-prob]{Teorema}
\newtheorem{thm:integral de lebesgue}[thm:base-prob]{Teorema}
\newtheorem{thm:convergencia monotona}[thm:base-prob]{Teorema}
\newtheorem{thm:convergencia dominada}[thm:base-prob]{Teorema}
\newtheorem{thm:toneli}[thm:base-prob]{Teorema}
\newtheorem{thm:fubini}[thm:base-prob]{Teorema}
\newtheorem{thm:integrcao contagem}[thm:base-prob]{Teorema}
\newtheorem{thm:holder}[thm:base-prob]{Teorema}
\newtheorem{thm:minkowski}[thm:base-prob]{Teorema}
\newtheorem{thm:consq holder}[thm:base-prob]{Teorema}
\newtheorem{thm:radon nikodym}[thm:base-prob]{Teorema}
\newtheorem{thm:radon nikodym2}[thm:base-prob]{Teorema}
\newtheorem{thm:tfc qs}[thm:base-prob]{Teorema}
\newtheorem{thm:mudanca de variavel}[thm:base-prob]{Teorema}
\newtheorem{thm:soma prob}[thm:base-prob]{Teorema}
\newtheorem{thm:probabilidade total}[thm:base-prob]{Teorema}
\newtheorem{thm:bayes}[thm:base-prob]{Teorema}

% Corolário
\newtheorem{cor:base-prob}{Corolário}[chapter]
\newtheorem{cor:equiprovavel}[cor:base-prob]{Corolário}
\newtheorem{cor:formula de bayes}[cor:base-prob]{Corolário}
